% ===========================================================================
% chapters/04_results.tex
% Updated after E1 multi-model results: Pythia-70M, 160M, 410M, 1B
% ===========================================================================
\chapter{Evaluation: Results and Discussion}
\label{ch:results}

This chapter reports the first completed observational experiment, E1, and then states the
remaining experiments in the order in which they will be used to test the stronger claims of the
thesis. E1 addresses the phase-identification premise: whether the weight space of a language
model passes through a reproducible early reorganisation during pre-training. The later experiments
ask whether that reorganisation is functionally meaningful (E2) and whether it marks a period of
heightened durable plasticity under intervention (E3--E6). Thus the evidence is deliberately
layered: E1 locates a candidate developmental window, E2 asks whether it has behavioural
correlates, and E3 tests the causal sensitive-period claim.

The current results establish the E1 observation across four Pythia-deduped model sizes. The
remaining sections are therefore still partly prospective, but the chapter is no longer only a
protocol: it now contains the first substantive result of the thesis.

\section{E1 finding: a reproducible early weight-space reorganisation}
\label{sec:res-prelim}

Tracking the weight-spectral indicators of \cref{sec:meth-indicators} across the Pythia checkpoint
schedule reveals a reproducible early reorganisation of weight space. The result appears in
Pythia-70M, 160M, 410M, and 1B, and is concentrated between roughly 128 and 2000 optimisation steps,
with the largest changes usually occurring between 512 and 2000 steps. This pattern is visible most
clearly in stable-rank collapse, spectral concentration, Marchenko--Pastur outlier growth in
rectangular matrices, changes in heavy-tail estimates, and instability of the leading singular
subspaces.

This result should be read carefully. It does not imply that the model has a single universal phase
boundary at one checkpoint. Rather, E1 identifies a short early developmental window during which
many modules rapidly depart from their near-initial spectral state. The precise peak interval shifts
with model size and module. Pythia-70M has its strongest broad stable-rank change between 1000 and
2000 steps, whereas Pythia-160M, 410M, and 1B peak most strongly between 512 and 1000 steps, with
continued reorganisation into 1000--2000 steps. The transition is therefore better described as a
window than as a point.

\begin{figure}[t]
  \centering
  \IfFileExists{figures/e1_multimodel_stable_rank_relative_early_8000.png}{%
    \includegraphics[width=0.92\linewidth]{figures/e1_multimodel_stable_rank_relative_early_8000.png}%
  }{%
    \fbox{\parbox[c][0.28\textheight][c]{0.88\linewidth}{\centering\itshape
    Figure placeholder: relative stable-rank trajectories for Pythia-70M, 160M, 410M, and 1B up to
    step 8000. Source artifact: \texttt{e1\_multimodel\_stable\_rank\_relative\_early\_8000.png}.}}%
  }
  \caption{Relative stable-rank trajectories during the early checkpoint window. Across four Pythia
  model sizes, most modules remain close to their initial spectral dimensionality at the earliest
  checkpoints and then undergo a sharp collapse before or around step 2000.}
  \label{fig:e1-stable-rank-relative}
\end{figure}

\begin{figure}[t]
  \centering
  \IfFileExists{figures/e1_multimodel_interval_mean_stable_rank_drop.png}{%
    \includegraphics[width=0.82\linewidth]{figures/e1_multimodel_interval_mean_stable_rank_drop.png}%
  }{%
    \fbox{\parbox[c][0.24\textheight][c]{0.82\linewidth}{\centering\itshape
    Figure placeholder: mean adjacent-checkpoint stable-rank drop by model and interval. Source
    artifact: \texttt{e1\_multimodel\_interval\_mean\_stable\_rank\_drop.png}.}}%
  }
  \caption{Adjacent-checkpoint stable-rank change by model size. The largest mean drops are
  concentrated in the intervals 128--512, 512--1000, and 1000--2000, with the strongest interval
  shifting slightly with model size.}
  \label{fig:e1-interval-drop}
\end{figure}

The module-level structure is as important as the global timing. Attention and MLP modules do not
move in lockstep. The attention query/key/value matrices tend to develop strong outlier structure
early. The MLP expansion projection, \texttt{mlp.dense\_h\_to\_4h}, often remains close to its
initial stable rank through the earliest checkpoints and then collapses sharply between 512 and
2000 steps. In the 1B model, for example, its mean stable rank is approximately 913.6 at step 0,
913.4 at step 128, and 853.3 at step 512, before dropping to about 209.5 at step 1000 and 59.7 at
step 2000. This pattern argues against a simple monotone-drift interpretation: at least some large
weight matrices remain near-initial for several released checkpoints before undergoing a rapid
spectral reorganisation.

\begin{table}[t]
  \centering
  \small
  \begin{tabular}{lrrrr}
    \toprule
    Model & Matrix--checkpoint rows & Layers & Modules & Checkpoints \\
    \midrule
    Pythia-70M  & 240 & 6  & 4 & 10 \\
    Pythia-160M & 480 & 12 & 4 & 10 \\
    Pythia-410M & 960 & 24 & 4 & 10 \\
    Pythia-1B   & 640 & 16 & 4 & 10 \\
    \midrule
    Total & 2320 & -- & -- & -- \\
    \bottomrule
  \end{tabular}
  \caption{Coverage of the E1 multi-model spectral analysis. Each row corresponds to one module
  matrix at one checkpoint. The checkpoint grid is \texttt{step0}, \texttt{step128},
  \texttt{step512}, \texttt{step1000}, \texttt{step2000}, \texttt{step3000}, \texttt{step8000},
  \texttt{step16000}, \texttt{step64000}, and \texttt{step143000}.}
  \label{tab:e1-coverage}
\end{table}

\begin{table}[t]
  \centering
  \small
  \begin{tabular}{lccc}
    \toprule
    Model & Strongest interval & Mean stable-rank drop & Modules $>10\%$ drop \\
    \midrule
    Pythia-70M  & 1000--2000 & 43.1\% & 4/4 \\
    Pythia-160M & 512--1000  & 49.3\% & 4/4 \\
    Pythia-410M & 512--1000  & 51.0\% & 4/4 \\
    Pythia-1B   & 512--1000  & 58.9\% & 4/4 \\
    \bottomrule
  \end{tabular}
  \caption{Strongest adjacent-checkpoint stable-rank drops by model size. The exact peak interval
  varies, but the strongest changes consistently occur inside the early 512--2000-step window.}
  \label{tab:e1-peak-intervals}
\end{table}

Two caveats bound the interpretation. First, these are weight-space phases, not direct claims about
activation geometry. They may correspond to representation-space transitions, but that
correspondence is not assumed here. Second, E1 is observational. It identifies a candidate
developmental window; it does not by itself establish that the model is more durably plastic during
that window. That stronger statement requires the behavioural and intervention experiments that
follow.

\section{E1 --- Phase identification and validation}
\label{sec:res-e1}

\noindent\textbf{Hypothesis.} \Cref{hyp:phases}. \quad\textbf{Status.} Completed as an
observational experiment across Pythia-70M, 160M, 410M, and 1B; causal interpretation deferred to
E3.

\noindent\textbf{Protocol.} E1 computes the full indicator panel over the Pythia checkpoint grid,
per layer and per module, for \texttt{attention.query\_key\_value}, \texttt{attention.dense},
\texttt{mlp.dense\_h\_to\_4h}, and \texttt{mlp.dense\_4h\_to\_h}. The analysis uses both absolute
and relative trajectories, adjacent-checkpoint change scores, linear-step plots, and log-step plots.
The checkpoint grid is fixed across the four model sizes: step 0, 128, 512, 1000, 2000, 3000, 8000,
16000, 64000, and 143000.

\noindent\textbf{Result.} The phase-identification hypothesis is supported at the level of
weight-space geometry. Across all four model sizes, the largest stable-rank drops concentrate before
step 2000, especially in the 512--1000 and 1000--2000 intervals. The result survives the most
important visual validation: it is visible on the early-window linear axis and is not merely a
log-axis artefact. It also replicates across size: the exact peak interval changes, but the early
window itself is stable.

The clearest primary signal is stable-rank collapse. Effective rank moves less sharply, especially
at larger scale, and should therefore be treated as a weaker corroborating measure rather than as
the main boundary detector. MP-outlier counts are informative for the rectangular QKV and MLP
matrices, but not equally useful for every module shape. This is why the boundary is credited only
from convergence across several indicators and not from any one spectral statistic.

\noindent\textbf{Interpretation.} E1 supports the existence of a candidate developmental window
rather than a single global boundary. It also supports the module-specific direction of the thesis:
attention and MLP matrices reorganise at different times, and the MLP expansion projection shows a
particularly clear delayed collapse. The candidate checkpoint grid for later intervention is
therefore
\[
  \{0, 128, 512, 1000, 2000, 3000, 8000\},
\]
which covers the near-initial period, the rising edge of reorganisation, the strongest transition
intervals, and a post-window control.

\noindent\textbf{Claim supported.} E1 establishes a replicated observational result: early
pre-training contains a short, module-staggered weight-space reorganisation across Pythia model
sizes from 70M to 1B. It does not establish a sensitive or critical period. It supplies the
independent geometric window that later experiments must test behaviourally.

\section{E2 --- Functional grounding}
\label{sec:res-e2}

\noindent\textbf{Hypothesis.} \Cref{hyp:functional}. \quad\textbf{Status.} Planned next.

\noindent\textbf{Protocol.} Align the E1 window with loss and lightweight behavioural probes over
the identical checkpoints (\cref{sec:meth-functional}). The first pass uses fixed-text next-token
loss and simple local probes, rather than making the result depend on large external evaluation
suites.

\noindent\textbf{Discriminating prediction / outcomes.} If loss curvature or probe performance
changes in the same early window, then the weight-space transition is functionally grounded. If the
curves are smooth or unrelated to the E1 window, the E1 phases remain valid as geometric phases but
are not yet functionally consequential.

\resultplaceholder{Overlay of E1 phase windows on loss and per-probe trajectories; quantitative test
of boundary--acquisition coincidence.}

\section{E3 --- The critical-period intervention (causal core)}
\label{sec:res-e3}

\noindent\textbf{Hypothesis.} \Cref{hyp:critical}. \quad\textbf{Status.} Planned after E2.

\noindent\textbf{Protocol.} Inject a defined signal at checkpoints spanning the E1 window; measure
uptake, retention after fixed subsequent training, and degradation-resistance
(\cref{sec:meth-intervention}); plot durability against injection step and overlay the E1 window.

\noindent\textbf{Discriminating prediction.} A genuine sensitive period predicts a break in
durability aligned with the E1 consolidation boundary, exceeding the learning-rate,
cumulative-exposure, and monotone ``earlier is better'' baselines (\cref{sec:meth-justification}).

\noindent\textbf{Pre-registered outcomes.} A boundary-aligned durability break, beating baselines,
would support the critical-period hypothesis. A smooth monotone decay would be an honest negative:
it would show that early training matters, but not that there is a period. No dependence on
injection time would count against the sensitive-period interpretation.

\resultplaceholder{Durability-vs-injection-step curves for uptake, retention, and
regradation-resistance, with the E1 window overlaid; baseline comparisons; confidence intervals over
seeds and checkpoints.}

\section{E4 --- Generality across signal types}
\label{sec:res-e4}

\noindent\textbf{Hypothesis.} Generality clause of \cref{hyp:structure}. \quad\textbf{Status.}
Planned.

\noindent\textbf{Protocol.} Repeat E3 with two or three distinct signal types, such as a synthetic
factual signal, an algorithmic or capability signal, and an alignment-style proxy; compare window
locations.

\noindent\textbf{Discriminating prediction / outcomes.} A shared window across signal types would
support the interpretation that sensitivity is a property of developmental stage. Signal-specific
windows would imply a weaker but still informative conclusion: that the timing of durable
plasticity depends on the kind of signal.

\resultplaceholder{Durability-vs-step curves per signal type on common axes; test of shared versus
distinct boundaries.}

\section{E5 --- Module-specific plasticity}
\label{sec:res-e5}

\noindent\textbf{Hypothesis.} Module-specificity clause of \cref{hyp:structure}.
\quad\textbf{Status.} Planned.

\noindent\textbf{Protocol.} The first version treats E5 as an analysis attached to E3 rather than as
an expensive separate intervention. During injection, per-module gradient and update norms are
recorded. The analysis then asks whether retention is better predicted by the consolidation timing
of the most-affected module than by the global E1 window.

\noindent\textbf{Discriminating prediction / outcomes.} If durability tracks each component's own
consolidation timing, this would support a module-resolved sensitive period. If no module dependence
appears, the data support a single global window or no window at all.

\resultplaceholder{Per-module update norms, per-module consolidation times, and regression relating
module-specific consolidation to durability.}

\section{E6 --- Alignment-style durability case}
\label{sec:res-e6}

\noindent\textbf{Hypothesis.} Application clause of \cref{hyp:structure}. \quad\textbf{Status.}
Planned as a scoped application, not as a load-bearing safety claim.

\noindent\textbf{Protocol.} Inject an alignment-style proxy signal across the E1 window and compare
its retention with late/post-hoc injection. At this scale, the degradation test is framed as
conflicting or adversarially directed fine-tuning of the proxy behaviour, not as a claim about
frontier-model safety.

\noindent\textbf{Discriminating prediction / outcomes.} If in-window alignment-style signal is more
retained after degradation than late signal, the developmental result has a safety-relevant
application. If not, the developmental E3 result can still stand without the alignment application.

\resultplaceholder{Proxy robustness metrics for in-window versus post-hoc alignment-style signal,
with explicit small-scale caveats.}

\section{Toy-model calibration and mechanism}
\label{sec:res-toy}

\noindent\textbf{T1 (calibration).} T1 is deferred. The first pilot was useful as an implementation
check but did not produce the required behavioural generalisation transition: the toy model
memorised the training set while held-out accuracy did not improve. Consequently, the intervention
retention readout was not interpretable. T1 therefore cannot currently be used as calibration
evidence for the indicator-to-transition interpretation.

\noindent\textbf{T2 (mechanistic vignette).} A mechanistic toy vignette remains optional. It should
only be included if a revised toy task first produces a known transition that the indicators can
recover.

\resultplaceholder{Revised toy calibration results, if run: indicator panel across a known
transition; inject-before/after durability; and any sufficiency demonstration.}

\section{Summary of evidence}
\label{sec:res-summary}

The current evidence supports \cref{hyp:phases} at the observational, weight-space level. Across
Pythia-70M, 160M, 410M, and 1B, E1 identifies a reproducible early reorganisation concentrated
before step 2000 and strongest between 512 and 2000 steps. The result is module-staggered rather
than globally synchronous. This is sufficient to define the checkpoint window for the next
experiments, but it is not yet sufficient to claim a sensitive or critical period. The next required
step is E2 functional grounding, followed by E3 intervention.
